\chapter{视觉 Transformer}
\label{lab:E40_cv_vit}

\noindent\textbf{配套文件：}\path{notebook/E40_cv_vit.ipynb}\par
\noindent\textbf{教材关联：}\bookxrefshort{ch:vision-transformer}。

\section*{实验目标}
验证patch排列和视觉编码器接口。

\section*{准备及定位}
先阅读本册实验总则，确认Notebook中的依赖与资产单元。原理计算、库接口迁移和真实模型训练可能具有不同资源要求，应分别记录设备、版本及运行范围。

源码定位可从下列函数开始；应结合前置数据与配置单元阅读。
\begin{itemize}
\item \texttt{my\_patchify}
\end{itemize}

\section*{操作及观察}
\begin{enumerate}
\item 使用可追踪的图像数组检查patchify的像素顺序。
\item 比较补丁切分加线性投影与对应卷积投影的结果。
\item 核对CLS、位置编码及编码器输出形状，检查输入分辨率与补丁数。
\item 对照torchvision或Transformers接口，迁移分类头时记录预处理和权重来源。
\end{enumerate}

\section*{结果判读及提交}
提交补丁可视化、等价误差和预处理配置。形状一致不能排除像素顺序错误；改变分辨率需重新处理位置表示。

提交时附上所执行单元范围、环境与制品身份、原始结果及失败说明。将未运行项目明确列出，不能用Notebook中已有输出替代本次执行证据。

相关方法的原始定义与适用前提见\cite{dosovitskiyImageWorth16x162022}；本实验的运行结果仍须按实际配置单独验证。
